\section{Strategic Opacity}\label{sec:endogenousopacity}

This section analyzes endogenous opacity by allowing the manager to control $\omega_{e}$. It uncovers the manager's incentive to intentionally make the fund opaque and how it varies depending on her expected skill.

\subsection{Setup} \label{sec:setupEndog}
Prior to investors' capital allocation and the manager's investment decision, the fund manager chooses the fund opacity $\omega_{e}$, anticipating the equilibrium reactions of investors and the market maker specified in Sections \ref{sec:eqm} and \ref{sec:opaque fund}. The manager seeks to maximize her \textit{ex-ante} expected utility derived from the fee income, that is, $U\equiv \mathrm{E}[\phi y]$. By incorporating the fund flow in equation \eqref{eq:M}, her objective function is specified as follows.\footnote{The law of iterated expectation implies $\mathrm{E}\left[\mathrm{E}[\pi|f]\right] = \mathrm{E}[\pi]$.}
\begin{equation}
U = \phi \mathrm{E}[\pi + M] = \phi(1+\theta)\mathrm{E}[\pi] - c(\omega_e).\label{eq:U}
\end{equation}
To formalize the manager's optimization problem in the rational expectations framework, we assume that, when parameter values support multiple equilibria, all agents in the economy coordinate their beliefs toward the low-$\beta$ equilibrium.\footnote{Alternatively, we may assume that the agents rely on the realization of a sunspot shock to coordinate their beliefs, where they anticipate the low- and the high-$\beta$ equilibria with probability $\rho\in[0,1]$ and $1-\rho$, respectively. Our setting corresponds to $\rho=1$, and this assumption does not affect our qualitative results.} Also, for the sake of tractability, we assume that $c'(\omega_e)$ is sufficiently large, so that the manager's maximization problem satisfies the second-order condition.

In this extended model, the definition of equilibrium augments Definition \ref{def:eqm} of the baseline model by additionally requiring that $\omega_e$ maximizes the manager’s objective function $U$ in equation \eqref{eq:U}.

% introduce an equilibrium-selection device in cases of multiple equilibria. In particular, following the most common approach, we assume that market participants condition their actions on the realization of an exogenous \textit{sunspot} shock, characterized by the binary random variable $z\in\{0,1\}$.\footnote{The sunspot equilibrium is proposed by \citet{diamond1983bank} as an equilibrium selection device in the context of bunk runs and formalized in the equilibrium analyses by the subsequent studies, such as \citet{cooper1998bank}.} Namely, whenever parameters admit multiple equilibria, market participants play the high-$\beta$ equilibrium if a spot appears on the sun ($z=1$), whereas the low-$\beta$ equilibrium is realized if no spots are observed ($z=0$). Accordingly, the high-$\beta$ equilibrium is realized with probability $\rho_{H}\equiv\Pr(z=1)$, while the complementary probability is denoted as $\rho_{L} = 1-\rho_{H}$. We assume that $z$ is realized at the beginning of the trading stage in $t=1$. In what follows, the trading intensity in the high-$\beta$ and low-$\beta$ equilibria are denoted as $\beta = \beta_{H}$ and $\beta_{L}$, respectively. Also,  for consistency, we define $\rho_{j} = 1$ for $j=H,L$ when the equilibrium with $\beta = \beta_{j}$ is unique.
\subsection{Opacity as Commitment Device}
Given the trading intensity $\beta$, the unconditional expected trading profit in \eqref{eq:U} is computed as 
\begin{equation}
     \mathrm{E}[\pi] = \frac{\beta \omega_{u}}{\beta^2 \omega_u + \omega_s}\omega_{s}.\label{eq:unconditional pi}
\end{equation}
As in the standard \citet{kyle1985continuous} framework, the first fraction represents the profit margin for each unit of the informational advantage of the manager, and the second term $\omega_s$ represents the expected size of this advantage.

The fund opacity does not directly show up in equation \eqref{eq:unconditional pi}, but it influences the expected performance through the equilibrium trading intensity $\beta$.
\begin{corollary}
     The expected fund performance, defined by $\mathrm{E}[\pi]$, is monotonically increasing in the fund opacity.
 \end{corollary}
 \input{Figures/fig_opacity_beta_profit}
 Figure \ref{fig:opacity_pi} plots $\mathrm{E}[\pi]$ in relation to the fund opacity $\omega_{e}$, which traces the reactions of $\beta$ in Figure \ref{fig:opacitybeta}. In our model, $\mathrm{E}[\pi]$ declines with the trading intensity $\beta$. This is due to the presence of the showing-off motive: althrough the unconditional expected performance is maximized at the \citet{kyle1985continuous}'s benchmark with $\theta=0$, the manager seeks to trade more aggressively to attract fund flows. It revelas too much private information to the market maker and lowers the fund performance, while the manager is compensated for this loss by attracting more fund flows. The opacity, conversely, diminishes the showing-off motive and induces the manager to take more profit-driven strategies, bringing $\beta$ closer to the profit-maximizing level. As a consequence, opacity helps improve the expected fund performance and, according to the manager's utility \eqref{eq:U}, this effect encourages the manager to make her fund opaque.
 
 However, given that fund flow $M$ increases with the expected performance (see equation [\ref{eq:M}]), the showing-off motive may seem paradoxical: why does the manager trade aggressively to begin with, knowing it harms performance and reduces $M$? This tension arises from a crucial information timing gap. When she decides on the trading strategy $x$, the manager cannot alter the market’s beliefs about her strategy ($\beta$ and $\lambda$ in equations [\ref{eq:cond_pi}]--[\ref{eq:cond_var}]), but she can influence investor perceptions via the signal $f$ by inflating her trading volume $|x|$ to appear more skilled, creating the showing-off motive. Conversely, at the \textit{ex-ante} stage (i.e., before she sends the signal $f$ to investors), the manager recognizes that this future showing-off motive will lead to flashy but inefficient risk-taking that erodes unconditional performance and the expected fund flow $\mathrm{E}[M]$. Consequently, she strategically chooses fund opacity as a commitment device to ``tie her own hands," thereby dampening the investors' future responsiveness to $f$ and forcing her future self to prioritize profit-driven trading over reputation-driven trading.
 
 %Once again, multiple equilibria support somewhat non-trivial impacts of fund opacity to the fund's performance. Namely, given a relatively high fund-performance sensitivity (panel [b]), the market's belief about moderately opaque funds can coordinate on substantially different equilibria, leading similar funds to exhibit divergent performance outcomes. For example, two dots on panel (b) represent funds with the same opacity level, but they substantially differ in terms of the expected performance, depending on whether the market believes the realization of the high-$\beta$ (the lower dot) or the low-$\beta$ equilibrium (the upper dot). This result implies that empirical studies could find weak or even contradictory correlations between fund transparency and returns.

\subsection{Implications}
\subsubsection{Skill, Opacity, and Fragility}
We first propose ``middling" skill acquisition due to the multiple trading-stage equilibria by focusing on $\omega_{s}\in (\ubar{\omega}_{s},\bar{\omega}_{s})$. We refer to it as middling because the acquired skill level is either higher or lower than what she would acquire if she knew which type of equilibrium would be realized. On the equilibrium path, both the selected skill potential ($\omega_{s}^{*}$) and the market's belief about it ($\omega_{s}$) fall in this region if the marginal skill-acquisition cost is intermediate. In this region, the trading stage involves multiple equilibria (i.e., $\beta_{0} = \beta_{H}$ or $\beta_{L}$), whereas the trader at the skill-acquisition stage cannot influence the equilibrium selection.
%\footnote{The equilibrium realization in the trading stage is governed by the market's belief about the skill potential, $\omega_{s}$, as illustrated in Proposition \ref{prop:endogenous}, whereas the trader cannot influence it by selecting $\omega_{s}^{*}$ due to the lack of a commitment device.} 
Consequently, she determines her skill potential by anticipating the sunspot shock, which leads to the skill acquisition as follows. 
%We refer to it as \textit{conservative} because the acquired skill potential becomes either too high or too low compared to the skill potential she would acquire if she were allowed to select it again after the trading-stage equilibrium is realized.

%\begin{proposition}[Middling skill acquisition]\label{prop:skill_inef}
%The equilibrium skill potential in the presence of multiple trading-stage equilibria is middling in the sense that; 
%\begin{enumerate}[label=(\roman*)]
%\item It is lower than the skill potential that she would acquire in the low-$\beta$ equilibrium; and 
%\item It is higher than the skill potential that she would acquire in the high-$\beta$ equilibrium.
%\end{enumerate}
%\end{proposition}

\begin{proposition}[Middling skill acquisition]\label{prop:skill_inef}
The equilibrium skill potential in the presence of multiple trading-stage equilibria is
\begin{enumerate}[label=(\roman*)]
\item lower than the skill potential that she would acquire in the low-$\beta$ equilibrium; and 
\item higher than the skill potential that she would acquire in the high-$\beta$ equilibrium.
\end{enumerate}
\end{proposition}

Figure \ref{fig:skill_potential} illustrates a numerical example; $\omega_{s,m}$ is the ex-ante optimal skill potential, while the ex-post equilibrium is either $\omega_{s,\ell}$ or $\omega_{s,h}$, depending on the equilibrium
realization in the trading stage.

The trader invests in skill potential incorporating the possibility of both high- and low-$\beta$ equilibria. At the trading 
stage, however, the acquired skill potential turns out to be insufficient if the low-$\beta$ equilibrium is realized, as its optimal play involves the profit-driven strategy and requires a relatively high asset-selection skill (i.e., $\omega_{s,h}$ is the \textit{ex-post} equilibrium). Conversely, the acquired skill potential is too high if the high-$\beta$ equilibrium is realized, as the trader takes the reputation-driven strategy that does not require a high skill potential (i.e., $\omega_{s,\ell}$ is the \textit{ex-post} equilibrium). 

This result indicates that individuals or firms may invest in skills that seem oddly intermediate or cautious—neither fully specialized nor overly basic. For example, it explains why some traders and finance professionals choose to pursue broader business qualifications instead of more specialized academic paths, allowing them to adapt to a wide range of market situations. Traders also learn intermediate-level coding (e.g., Python) for algorithmic trading or data analysis but stop short of deep computer science expertise. A similar trend can be observed with respect to CFA Levels and quantitative skills among hedge fund managers. These skill acquisitions would be difficult to rationalize in an economy with a single equilibrium but become logical under the equilibrium multiplicity and uncertainty regarding equilibrium realization. 

\subsubsection{Fragile Skill Acquisition and Skill Diversity}
By expanding our scope to the entire range of skill potential ($\omega_{s}$), the discontinuities in Lemma \ref{lem:MB_beta} imply that the skill-acquisition problem in (\ref{eq:skillacquisition}) may have multiple solutions.\footnote{Although the equilibrium multiplicity in the trading stage (i.e., $\beta_{0}$) is determined by the level of $\omega_{s}$, as Proposition \ref{prop:endogenous} stipulates, it is orthogonal to the existence of multiple solutions to the optimal skill acquisition in equation (\ref{eq:FOC_SA}).} 

For example, when the marginal cost of skill acquisition is moderately small, as Figure \ref{fig:skill_potential} illustrates, both the high skill potential ($\omega_{s,h}$) and the intermediate skill potential ($\omega_{s,m}$) become optimal for the trader. Corresponding to $\omega_{s,h}$, the equilibrium in the trading stage is unique with $\beta_{0}=\beta_{L}$, where the trader engages in solid profit-driven trading. Conversely, the trading stage with skill potential $\omega_{s,m}$ involves multiple equilibria, and the selected skill potential is middling, as Proposition \ref{prop:skill_inef} demonstrates. The symmetric argument holds when the marginal cost is relatively high: the intermediate skill potential and the low skill potential show up as multiple optima. This result can explain at least three realistic phenomena without relying on ad-hoc assumptions on the skill-acquisition environment.

Firstly, the coexistence of traders with different skill levels arises not from heterogeneity in skill-acquisition problems (e.g., different marginal costs), but from the multiple trading-stage equilibria and the possibility of mixed strategies in the skill acquisition stage. This reflects a realistic scenario where, due to advancements in the internet and social media (e.g., online webinars via Zoom and academic courses uploaded on YouTube), opportunities for skill acquisition have become widely and uniformly accessible to all traders, reducing barriers to learning. Despite the increased homogeneity in skill-acquisition costs, we continue to observe heterogeneity in skill levels in the finance industry. Our model offers an explanation for this observation, showing how diverse skill levels can emerge among market participants, without the need to assume inherent differences in their skill acquisition costs.

Secondly, it also captures the possibility of widespread surges in skill acquisition without assuming special conditions for the skill-acquisition environment. For example, we can explain traders and fund managers in competitive environments (e.g., Wall Street) may collectively strive to acquire high-level skills, such as pursuing advanced degrees, driven by the mixed-strategy equilibrium. This surge in skill acquisition could occur even in the absence of significant changes in external factors, and without the need for heterogeneity in skill-acquisition costs, reflecting industry-wide trends in upskilling and degree inflation, as documented by \citet{philippon2012wages} and \citet{fuller2017dismiss}.

Finally, discontinuities in the marginal benefit curve add another layer of insight into the effect of the skill-acquisition cost on the equilibrium skill potential. Namely, discontinuities introduce the possibility of significant jumps in the optimal skill level triggered by a small change in the marginal skill-acquisition cost. For instance, when marginal technological advancements slightly reduce the marginal cost of learning, the result could be a sizable upward shift in skill acquisition. Conversely, even a small increase in the cost of learning might cause a large decline in the equilibrium skill level. This discontinuity-driven sensitivity can help explain the historical observations stylized by \citet{philippon2012wages} that certain external changes lead to abrupt shifts in the distribution of skill levels in the finance industry.